% ======================================================================
% ABSTRACT (target 150–180 words)
% Order: problem -> gap -> design -> local performance -> full-hand
% performance -> airflow/pulse -> implication.
% Every bracketed value is a hard placeholder; do not soften or remove
% the \dataneeded marks until real numbers exist.
% ======================================================================
\section*{Abstract}

Robotic and wearable systems require compliant tactile interfaces that can
resolve weak, rapidly varying and spatially distributed mechanical stimuli.
Existing tactile skins commonly trade off weak-force sensitivity, multi-axis
perception, spatial coverage and readout complexity. Here we report a magnetic
cilia tactile skin that combines a dense passive layer of magnetized elastomer
cilia with distributed tri-axis magnetic readout. Geometry optimization across
six cilia designs establishes \dataneeded{optimal length $\times$ diameter,
e.g.\ $L$\,=\,x\,mm, $d$\,=\,x\,mm} and yields a sensitivity enhancement of
\dataneeded{gain $G_{\mathrm{cilia}}$ with 95\% CI, $n \geq 3$ patches}-fold
relative to a matched flat magnetic control. A local $2\times2$ sensing unit
reconstructs unseen contact positions with a localization error of
\dataneeded{RMSE in mm, unseen-position test set} and estimates normal force
with an error of \dataneeded{NRMSE, \%}. Integrated into a 44-node hand-shaped
array, the system generates distributed tactile maps at \dataneeded{measured
full-frame rate, Hz} while suppressing motion-induced magnetic artifacts by
\dataneeded{CMR, dB}. The same platform resolves airflow between
\dataneeded{min wind speed, m\,s$^{-1}$} and \dataneeded{max wind speed,
m\,s$^{-1}$} and records radial-artery pulse micro-vibrations synchronized
with electrocardiography across \dataneeded{number of subjects}. These results
establish magnetic cilia as a scalable mechanical transduction layer for
embodied perception of subtle external and physiological stimuli.
